\documentclass[11pt]{article}

\usepackage[margin=1in]{geometry}
\usepackage{booktabs}
\usepackage{graphicx}
\usepackage[hidelinks]{hyperref}
\usepackage{xcolor}
\emergencystretch=2em

\title{SafeRL-Drive: Curriculum Learning for Procedural Driving in MetaDrive}
\author{Dheeraj Dhillon}
\date{\today}

% Phase-1 values are the exact fingerprint-compatible 20-episode test results.
\newcommand{\PpoSuccess}{100.0\%}
\newcommand{\SacSuccess}{100.0\%}
\newcommand{\IdmSuccess}{100.0\%}
\newcommand{\PpoCollision}{0.0\%}
\newcommand{\SacCollision}{0.0\%}
\newcommand{\IdmCollision}{0.0\%}
\newcommand{\PpoOffRoad}{0.0\%}
\newcommand{\SacOffRoad}{0.0\%}
\newcommand{\IdmOffRoad}{0.0\%}
\newcommand{\PpoCompletion}{96.1\%}
\newcommand{\SacCompletion}{95.7\%}
\newcommand{\IdmCompletion}{95.6\%}

\newcommand{\DirectSuccess}{12.0\%}
\newcommand{\CurriculumSuccess}{87.0\%}
\newcommand{\DirectCollision}{15.5\%}
\newcommand{\CurriculumCollision}{0.0\%}
\newcommand{\DirectOffRoad}{83.0\%}
\newcommand{\CurriculumOffRoad}{9.5\%}
\newcommand{\DirectTimeout}{5.0\%}
\newcommand{\CurriculumTimeout}{3.5\%}
\newcommand{\DirectCompletion}{47.0\%}
\newcommand{\CurriculumCompletion}{95.1\%}

\newif\ifphaseoneresults
\newif\ifphasetworesults
\IfFileExists{generated_phase1_results.tex}{
  \phaseoneresultstrue% Generated from the compatible 2026-07-22 Phase-1 comparison.
\renewcommand{\PpoSuccess}{100.0\%}
\renewcommand{\SacSuccess}{100.0\%}
\renewcommand{\IdmSuccess}{100.0\%}
\renewcommand{\PpoCollision}{0.0\%}
\renewcommand{\SacCollision}{0.0\%}
\renewcommand{\IdmCollision}{0.0\%}
\renewcommand{\PpoOffRoad}{0.0\%}
\renewcommand{\SacOffRoad}{0.0\%}
\renewcommand{\IdmOffRoad}{0.0\%}
\renewcommand{\PpoCompletion}{96.1\%}
\renewcommand{\SacCompletion}{95.7\%}
\renewcommand{\IdmCompletion}{95.6\%}

}{
  \IfFileExists{reports/generated_phase1_results.tex}{
    \phaseoneresultstrue% Generated from the compatible 2026-07-22 Phase-1 comparison.
\renewcommand{\PpoSuccess}{100.0\%}
\renewcommand{\SacSuccess}{100.0\%}
\renewcommand{\IdmSuccess}{100.0\%}
\renewcommand{\PpoCollision}{0.0\%}
\renewcommand{\SacCollision}{0.0\%}
\renewcommand{\IdmCollision}{0.0\%}
\renewcommand{\PpoOffRoad}{0.0\%}
\renewcommand{\SacOffRoad}{0.0\%}
\renewcommand{\IdmOffRoad}{0.0\%}
\renewcommand{\PpoCompletion}{96.1\%}
\renewcommand{\SacCompletion}{95.7\%}
\renewcommand{\IdmCompletion}{95.6\%}

  }{\phaseoneresultsfalse}
}
\IfFileExists{generated_phase2_results.tex}{
  \phasetworesultstrue\input{generated_phase2_results.tex}
}{
  \IfFileExists{reports/generated_phase2_results.tex}{
    \phasetworesultstrue\input{reports/generated_phase2_results.tex}
  }{\phasetworesultsfalse}
}

\begin{document}

\maketitle

\begin{abstract}
SafeRL-Drive is a reproducible student study of closed-loop reinforcement learning in
MetaDrive. Phase 1 established a trustworthy training and evaluation pipeline on a
traffic-free straight road. Frozen PPO and SAC policies each achieved 100\% success on 20
held-out episodes, although action analysis showed that PPO solved this simple task with a
constant open-loop command while SAC learned state-dependent continuous control. Phase 2
therefore focuses on SAC and asks whether a staged road-geometry curriculum improves
generalization to unseen three-block procedural maps under the same maximum budget as
direct training.
Across two complete, fingerprint-compatible seed pairs, curriculum SAC achieved 87\%
held-out success against 12\% for direct SAC, while raising route completion from 47.0\%
to 95.1\%. A zero-shot light-traffic test then exposed the remaining limitation:
curriculum SAC retained 69\% success but collided in 21\% of episodes. The project has
therefore answered its geometry-curriculum question; one bounded traffic-adaptation study
is proposed as the final portfolio extension.
Resolved configurations, strict experiment fingerprints, logs, per-episode metrics,
plots, videos, and Google Drive artifacts make each claim auditable.
\end{abstract}

\section{Introduction}
Closed-loop simulation evaluates a driving policy through the consequences of its own
actions. MetaDrive is particularly appropriate because its procedural generator can create
diverse road structures and disjoint scenario sets for generalization experiments
\cite{li2022metadrive}. The project began with PPO and SAC because they provide standard
on-policy and off-policy learning controls \cite{schulman2017ppo,haarnoja2018sac}. The aim
is not production autonomy. It is a compact, reproducible experiment that demonstrates
learning, exposes failure modes honestly, and tests one meaningful generalization
hypothesis.

\section{Scope and Success Criteria}
Phase 1 is a pipeline and basic-control gate. A frozen policy qualifies with at least 80\%
success and 90\% mean route completion, with no more than 10\% collision, off-road, or
timeout outcomes on 20 untouched episodes. Phase 2 changes one main dimension: road
geometry. Its target uses three procedurally generated blocks, full continuous control,
no training traffic, 100 training scenarios, 25 validation scenarios, and 100 disjoint
test scenarios. The same outcome thresholds are retained, and a train-to-test success gap
above 15 percentage points is treated as weak generalization.

Image observations, real-world scenario replay, multi-agent learning, vehicle-dynamics
randomization, obstacle curricula, and explicit safe-RL algorithms remain outside Phase 2.
Light traffic is used only as a zero-shot stress test after the traffic-free policy is
frozen.

\section{Phase-1 Setup}
Both learned agents used MetaDrive's bounded vector observation and reference reward.
Training used seed 0 on one straight topology with fixed spawning, no traffic, and a
500-step horizon. Validation used seeds 1000--1009. Final testing used seeds 4000--4019
and was not consulted during checkpoint selection. PPO used a $3\times3$ discrete
steering/throttle grid and CPU learning. SAC used continuous throttle/brake and steering
limited to $[-0.1,0.1]$ for the straight-road gate. Because these policy-facing action
interfaces differ, Phase 1 is a descriptive control matrix rather than a fair algorithm
ranking.

\section{Phase-1 Results}
\begin{table}[h]
  \centering
  \caption{Fingerprint-compatible held-out fixed-straight-road outcomes on seeds
  4000--4019.}
  \label{tab:phase1-results}
  \begin{tabular}{lcccc}
    \toprule
    Agent & Success & Collision & Off-road & Route completion \\
    \midrule
    IDM & \IdmSuccess & \IdmCollision & \IdmOffRoad & \IdmCompletion \\
    PPO & \PpoSuccess & \PpoCollision & \PpoOffRoad & \PpoCompletion \\
    SAC & \SacSuccess & \SacCollision & \SacOffRoad & \SacCompletion \\
    \bottomrule
  \end{tabular}
\end{table}

PPO achieved 100\% success, 96.1\% route completion, and zero reported failures on the
20 held-out episodes. Every deterministic action was exactly zero steering and full
throttle. Its sampled training rollouts remained noisy at the stopping checkpoint, with
3\% rolling success and 96\% off-road outcomes. The result proves that the discrete
pipeline can select a successful deterministic command, but not that PPO learned a robust
feedback controller.

SAC also achieved 100\% success and zero failures, with 95.7\% route completion. Its
sampled training success rose from 30\% at 30,000 steps to 58\%, 93\%, and 100\% by
45,000 steps. Held-out actions included small steering corrections, throttle modulation,
and occasional braking. SAC therefore supplies the better evidence of state-dependent
control and is the sole learned algorithm advanced to Phase 2.

For either 20-of-20 result, the 95\% Wilson interval for the underlying episode success
probability is 83.9--100\%. The point estimates pass the predeclared gate, but the interval
and single training seed rule out a strong robustness claim.

IDM achieved 100\% success, 95.6\% route completion, and zero failures on the same 20
held-out episodes. Its earlier validation reproducibility gate also produced identical
per-scenario outcomes across two fresh runs. The task fingerprints match for all three
agents. The result remains a descriptive task-outcome comparison because IDM, discrete
PPO, and continuous SAC expose different action interfaces.

\ifphaseoneresults
\begin{figure}[h]
  \centering
  \IfFileExists{../runs/phase1_comparison.png}{
    \includegraphics[width=0.92\linewidth]{../runs/phase1_comparison.png}
  }{
    \includegraphics[width=0.92\linewidth]{runs/phase1_comparison.png}
  }
  \caption{Fingerprint-checked Phase-1 held-out task outcomes. The controller interfaces
  differ, so the plot is descriptive rather than an algorithm ranking.}
\end{figure}
\fi

\section{Phase-2 Research Design}
Phase 2 tests whether a structured curriculum improves procedural-road learning. Both
conditions use Stable-Baselines3 SAC \cite{raffin2021sb3}, the same $[256,256]$ MLP,
full $[-1,1]$ steering and throttle/brake actions, identical rewards, a maximum of 500,000
environment steps, and the same final validation and test scenarios.

\paragraph{Direct SAC.}
The control is trained from scratch on 100 three-block procedural scenarios for the full
budget or until the target validation gate is reached.

\paragraph{Curriculum SAC.}
The experimental policy begins on curved roads (\texttt{C}, at most 100,000 steps), then
straight-plus-curve roads (\texttt{SC}, at most 150,000 steps), and finally uses all
remaining steps on random three-block maps. Each prerequisite stage must reach 90\%
success and route completion with at most 10\% collision, off-road, and timeout outcomes.
Model parameters and the SAC replay buffer persist across stages. Curriculum learning has
previously improved performance and sample efficiency across driving tasks of increasing
road complexity \cite{ozturk2021curriculum}; this experiment tests the idea in MetaDrive
under an equal total budget.

One direct and one curriculum seed-0 pilot are run first. Confirmation seeds 1 and 2 are
launched only if at least one pilot reaches 50\% success or 80\% route completion. A
curriculum advantage is claimed only for an improvement of at least ten percentage points
in success, or similar success with at least 20\% fewer combined collision and off-road
outcomes, with a consistent direction across the confirmation seeds.

\section{Phase-2 Results}
Before learned-policy training, both native controllers solved all 25 fixed validation
scenarios. IDM achieved 100\% success, 98.69\% mean route completion, and zero collision,
off-road, or timeout outcomes; ExpertPolicy achieved 100\%, 98.69\%, and zero respectively.
This establishes task feasibility, but it is not evidence that SAC has learned the task.

\begin{table}[h]
  \centering
  \caption{Unseen three-block procedural-road test results. Learned conditions are means
  across the two complete seed pairs (200 episodes per condition); native controllers use
  one 100-episode deterministic evaluation each.}
  \label{tab:phase2-results}
  \begin{tabular}{lccccc}
    \toprule
    Condition & Success & Collision & Off-road & Timeout & Route completion \\
    \midrule
    IDM & 95.0\% & 0.0\% & 5.0\% & 0.0\% & 95.9\% \\
    ExpertPolicy & 100.0\% & 0.0\% & 0.0\% & 0.0\% & 98.7\% \\
    Direct SAC & \DirectSuccess & \DirectCollision & \DirectOffRoad &
      \DirectTimeout & \DirectCompletion \\
    Curriculum SAC & \CurriculumSuccess & \CurriculumCollision &
      \CurriculumOffRoad & \CurriculumTimeout & \CurriculumCompletion \\
    \bottomrule
  \end{tabular}
\end{table}

Strict fingerprints matched across direct and curriculum conditions for algorithm,
architecture, action interface, reward, budget, map, and evaluation split. Curriculum
increased success by 75 percentage points, reduced off-road incidence by 73.5 points, and
increased mean route completion by 48.0 points. The direction was consistent in both
complete pairs: seed 0 improved from 24\% to 86\% success and seed 1 from 0\% to 88\%.
The curriculum aggregate passes the Phase-2 gate, although seed 0 alone retained a 14\%
off-road rate. Failure flags are not mutually exclusive, so collision, off-road, and
timeout percentages need not sum to 100\%.

The native policies remain useful ceilings rather than learned-policy equivalents. On the
same untouched 100 scenarios, IDM achieved 95\% success and ExpertPolicy 100\%. Curriculum
SAC is therefore close to IDM on traffic-free geometry, but it has not surpassed the
simulator's native controllers.

Optimization was materially seed-sensitive. Curriculum seed 2 failed the mandatory curve
prerequisite after 100,000 steps: its best validation checkpoint had 0\% success and
41.75\% route completion. Direct seed 2 completed its budget but had 0\% test success,
19.24\% route completion, and 100\% off-road incidence. These terminal results are kept
rather than rerun selectively. Thus the 87\% headline is conditional on the two complete
pairs, while curriculum produced a usable final procedural policy in two of three launches.

Checkpoint behavior also rejects the idea that more training alone is the answer.
Curriculum seed 0 passed the target gate after 150,000 total steps. Curriculum seed 1
reached 92\% validation success at 375,000 steps, then fell to 36\% at 500,000; the saved
best checkpoint still obtained 88\% on the untouched test. Direct seed 0 peaked at 40\%
validation success at 150,000 steps, whereas direct seeds 1 and 2 converged to almost
constant deterministic actions. Reliable checkpoint selection was therefore essential.

\begin{table}[h]
  \centering
  \caption{Zero-shot test at traffic density 0.05, averaged across the same two frozen
  learned-policy seeds (200 episodes per condition). Neither policy was trained with
  traffic.}
  \label{tab:phase2-traffic}
  \begin{tabular}{lccccc}
    \toprule
    Condition & Success & Collision & Off-road & Timeout & Route completion \\
    \midrule
    Direct SAC & 3.5\% & 37.0\% & 67.5\% & 4.0\% & 37.3\% \\
    Curriculum SAC & 69.0\% & 21.0\% & 6.0\% & 4.0\% & 87.1\% \\
    \bottomrule
  \end{tabular}
\end{table}

The light-traffic result is encouraging transfer, not traffic competence. Curriculum SAC
retained far more route-following ability than direct SAC, but a 21\% collision rate is
not a safety result. Seed outcomes also varied substantially: curriculum seeds 0 and 1
achieved 59\% and 79\% success, with 34\% and 8\% collision respectively. Traffic was a
deliberate distribution shift, and the default state observation already contains lidar
information about surrounding vehicles \cite{li2022metadrive}; the missing ingredient is
traffic-conditioned learning rather than a new image pipeline.

\begin{figure}[h]
  \centering
  \IfFileExists{../runs/phase2_comparison.png}{
    \includegraphics[width=0.92\linewidth]{../runs/phase2_comparison.png}
  }{
    \includegraphics[width=0.92\linewidth]{runs/phase2_comparison.png}
  }
  \caption{Strictly compatible direct-versus-curriculum SAC results on unseen procedural
  roads.}
\end{figure}

\section{Limitations}
Phase 1 used one training seed per algorithm and a task that admitted an open-loop
solution. Phase 2 remains state-vector based, traffic-free during training, and limited to
one simulator and two SAC training conditions. One hundred test episodes quantify
episode-level uncertainty, but two complete paired training seeds are too few for a broad
statistical claim. The failed third curriculum launch is a real availability limitation,
not missing data. Validation selected the best checkpoint, and the test set was not used
for that choice, but the aggregate is still conditional on completing the curriculum.
The zero-shot traffic test measures robustness but does not establish traffic competence
or safety. Finally, the two current top-down videos are technically valid 800-by-800
recordings of a fixed successful seed-30000 episode; they demonstrate behavior but are not
representative samples, especially for direct SAC.

\section{Bounded Final Extension}
The original MVP and the Phase-2 geometry question are complete. The only remaining
extension justified by the evidence is traffic adaptation: can a geometry-competent
curriculum policy be fine-tuned on modest traffic while retaining its traffic-free
generalization? This is a final demonstration, not a new algorithm search.

The experiment is deliberately capped at two long runs. The best curriculum checkpoints
from seeds 0 and 1 are each continued for at most 250,000 steps on the same three-block
procedural distribution, first at traffic density 0.02 and then at 0.05. Observation,
action interface, SAC architecture, reward, map complexity, and scenario split remain
fixed; only exposure to traffic changes. Direct-from-scratch training, new algorithms,
larger maps, image policies, and safe-RL objectives are excluded so that the result has one
interpretation.

Before training, IDM must establish that the exact held-out traffic tasks at densities
0.05 and 0.10 are solvable. Frozen pre- and post-adaptation policies are then compared on
the same 100 unseen scenarios at densities 0, 0.05, and 0.10. The primary density-0.05
gate, averaged across both seeds, is at least 80\% success and 90\% route completion with
no more than 10\% collision or off-road incidence. Traffic adaptation must also improve
success or collision by at least ten percentage points relative to the frozen zero-shot
policy, while traffic-free success may fall by no more than ten points. Density 0.10 is a
descriptive stress test, not a tuning target.

Reporting will include all two-seed outcomes whether favorable or not. The final visual
set will contain the existing traffic-free curriculum rollout and a matched traffic
failure/success pair on the same held-out scenario before and after adaptation. A
systematically selected failure case must be labeled as diagnostic rather than
representative. After these two training runs and frozen evaluations, the project stops;
SafeMetaDrive, multi-agent RL, real-data replay, and further seed or hyperparameter sweeps
remain future work.

\section{Conclusion}
Phase 1 converted an initially unreliable experiment into a reproducible control gate and
showed why headline success alone is insufficient: action behavior revealed that PPO's
perfect score was trivial, whereas SAC learned a responsive controller. Phase 2 returns to
the project's original procedural-generalization objective through one controlled
direct-versus-curriculum experiment. Its answer is positive but bounded: curriculum SAC
raised held-out success from 12\% to 87\% across the two complete pairs and approached the
IDM baseline on unseen traffic-free maps. The failed third curriculum launch and the 21\%
zero-shot traffic collision rate prevent a stronger reliability or safety claim. One
two-run traffic-adaptation extension is sufficient to turn that limitation into the
project's final experiment; beyond it, additional breadth would dilute rather than
strengthen the portfolio story.

\bibliographystyle{plain}
\bibliography{references}

\end{document}
